
\chapter{Pendahuluan}

\section{\label{sec:Accuracy}Akurasi}

\pagenumbering{arabic}
\fancyhead[LE,RO]{\thepage}
\fancyhead[RE]{\leftmark}
\fancyhead[LO]{\rightmark}

\subsection{Mengukur Galat}

Metode numerik dirancang untuk menghampiri berbagai macam objek. Kadang-kadang
yang dihampiri adalah akar, kadang-kadang turunan atau integral tentu,
kurva, ataupun solusi persamaan diferensial. Karena metode numerik
hanya menghasilkan hampiran terhadap objek-objek tersebut, penting
bagi kita untuk mengetahui seberapa akurat hasilnya. Terkadang akurasi
ditentukan oleh analisis aljabar yang cermat, terkadang oleh analisis
kalkulus yang cermat, dan sering kali oleh analisis polinomial Taylor
yang cermat. Namun, sebelum membahas perincian tersebut, kita perlu
membicarakan bagaimana galat dan, dengan demikian, akurasi diukur.

Ada dua ukuran dasar akurasi\index{akurasi}: galat absolut\index{galat!absolut}
dan galat relatif\index{galat!relatif}. Misalkan $p$ adalah nilai
yang hendak kita hampiri, dan $\tilde{p}$ merupakan hampiran terhadap
$p$. Maka $\tilde{p}$ meleset tepat sebesar $\tilde{p}-p$, yang
disebut galat\index{galat}. Tentu saja, $\tilde{p}-p$ akan bernilai
negatif ketika $\tilde{p}$ terlalu rendah. Artinya, ketika hampiran
$\tilde{p}$ lebih kecil daripada nilai eksak $p$. Sebaliknya, $\tilde{p}-p$
akan bernilai positif ketika $\tilde{p}$ terlalu tinggi. Namun, pada
umumnya kita tidak mempersoalkan apakah hampiran terlalu tinggi atau
terlalu rendah. Kita hanya ingin mengetahui seberapa besar penyimpangannya.
Karena itu, kita paling sering membicarakan galat absolut\index{galat!absolut},
$|\tilde{p}-p|$. Anda mungkin mengenali ekspresi $|\tilde{p}-p|$
sebagai jarak antara $\tilde{p}$ dan $p$, dan itu merupakan cara
yang baik untuk memandang galat absolut.

Galat absolut ketika menghampiri $p=\pi$ dengan bilangan rasional
$\tilde{p}=\frac{22}{7}$ adalah $|\frac{22}{7}-\pi|\approx0.00126$.
Galat absolut ketika menghampiri $\pi^{5}$ dengan bilangan rasional
$\frac{16525}{54}$ adalah $|\frac{16525}{54}-\pi^{5}|\approx0.00116$.
Galat absolut kedua hampiran ini hampir sama. Agar hal tersebut tampak
lebih jelas, $\pi\approx3.14159$ dan $\frac{22}{7}\approx3.14285$,
sedangkan $\pi^{5}\approx306.01968$ dan $\frac{16525}{54}\approx306.01851$.
Setiap hampiran mulai berbeda dari nilai eksaknya masing-masing pada
tempat perseribuan. Masing-masing hanya meleset sebesar 1 pada tempat
perseribuan tersebut.

Namun, ada hal lain yang terjadi di sini. $\pi$ mendekati 3, sedangkan
$\pi^{5}$ mendekati 300. Untuk menghampiri $\pi$ secara akurat hingga
perseratus terdekat, hampiran hanya perlu sama dengan nilai eksak
pada 3 nilai tempat---satuan, persepuluhan, dan perseratusan. Untuk
menghampiri $\pi^{5}$ secara akurat hingga perseratus terdekat, hampiran
harus sama dengan nilai eksak pada 5 nilai tempat---ratusan, puluhan,
satuan, persepuluhan, dan perseratusan. Dengan istilah yang lebih
ilmiah, kita mengatakan bahwa $\frac{22}{7}$ menghampiri $\pi$ secara
akurat hingga 3 angka signifikan, sedangkan $\frac{16525}{54}$ menghampiri
$\pi^{5}$ secara akurat hingga 5 angka signifikan. Di sinilah letak
inti galat relatif: membandingkan galat absolut dengan besar nilai
yang sedang dihampiri. Perbandingan ini dilakukan dengan menghitung
rasio galat terhadap nilai eksak. Dengan demikian, galat relatif ketika
menghampiri $\pi$ dengan $\frac{22}{7}$ adalah ${\displaystyle \frac{|\frac{22}{7}-\pi|}{|\pi|}\approx4.02(10)^{-4}}$,
sedangkan galat relatif ketika menghampiri $\pi^{5}$ dengan $\frac{16525}{54}$
adalah ${\displaystyle \frac{|\frac{16525}{54}-\pi^{5}|}{|\pi^{5}|}\approx3.81(10)^{-6}}$.
Kedua galat relatif berbeda dengan faktor sekitar 100 (setara dengan
kira-kira dua angka signifikan dalam akurasi), meskipun galat absolutnya
hampir sama. Secara umum, galat relatif\index{galat!relatif} ketika
menghampiri $p$ dengan $\tilde{p}$ diberikan oleh ${\displaystyle \frac{|\tilde{p}-p|}{|p|}}$.

\subsection{Sumber Galat}

Ada dua kategori umum galat. Galat algoritmik\index{galat!algoritmik}
dan galat titik-mengambang\index{galat!titik-mengambang}. Galat algoritmik
adalah setiap galat yang timbul dari metode hampiran itu sendiri.
Artinya, galat ini tidak dapat dihindari sekalipun kita melakukan
perhitungan eksak pada setiap langkah. Galat titik-mengambang timbul
karena komputer dan kalkulator pada umumnya tidak melakukan aritmetika
eksak, melainkan aritmetika titik-mengambang.

\noindent\begin{digression}{Standar IEEE 754}

\noindent Nilai titik-mengambang\index{aritmetika titik-mengambang}
disimpan dalam bentuk biner. Menurut Standar IEEE 754 yang digunakan
oleh kebanyakan komputer, mantisa (atau koefisien signifikan) disimpan
menggunakan 52 bit, yaitu 52 tempat biner. Karena bit terdepan selalu
dianggap bernilai 1 (dan karena itu tidak benar-benar disimpan), setiap
bilangan titik-mengambang direpresentasikan menggunakan 53 nilai tempat
biner yang berurutan. Sekarang mari kita perhatikan bagaimana $1/7$
direpresentasikan secara eksak. Dalam biner, sepertujuh sama dengan
$0.001001001\ldots$ karena $\frac{1}{7}=\sum_{i=1}^{\infty}2^{-3i}=\frac{1}{8}+\frac{1}{64}+\frac{1}{512}+\cdots$.
Untuk melihat bahwa hal ini benar, ingatlah dari kalkulus bahwa
\begin{eqnarray*}
\sum_{i=1}^{\infty}2^{-3i} & = & \sum_{i=1}^{\infty}\left(2^{-3}\right)^{i}\\
 & = & \frac{2^{-3}}{1-2^{-3}}\\
 & = & \frac{1/8}{7/8}\\
 & = & \frac{1}{7}.
\end{eqnarray*}
Namun, dalam Standar IEEE 754, $\frac{1}{7}$ dipangkas menjadi
\[
1.0010010010010010010010010010010010010010010010010010\times(2)^{-3}
\]
atau $\sum_{i=1}^{18}2^{-3i}$ yang tepat sama dengan $\frac{2573485501354569}{18014398509481984}$.
Dengan demikian, galat titik-mengambang dalam menghitung $1/7$ adalah
\begin{eqnarray*}
\left|\frac{2573485501354569}{18014398509481984}-\frac{1}{7}\right| & = & \frac{1}{126100789566373888}\approx7.93(10)^{-18}.
\end{eqnarray*}

\begin{description}
\item [{Referensi}] \cite{cunyieee754,goldberg}
\end{description}
\noindent\end{digression} 

Dalam aritmetika titik-mengambang, kalkulator atau komputer biasanya
menyimpan nilai dengan sekitar 16 angka signifikan. Sebagai contoh,
pada komputer atau kalkulator biasa (yang menggunakan aritmetika presisi
ganda), bilangan $\frac{1}{7}$ disimpan kira-kira sebagai $0.1428571428571428$,
sedangkan nilai eksaknya adalah $0.1428571428571428\ldots$. Pada
nilai eksak, pola $142857$ berulang tanpa henti, sedangkan pada nilai
titik-mengambang pengulangan berhenti setelah angka 8 yang ketiga.
Nilai tersebut dipangkas hingga 16 tempat desimal dalam representasi
titik-mengambang. Jadi, galat titik-mengambang dalam menghitung $1/7$
adalah sekitar $5(10)^{-17}$. Saya menggunakan kata ``sekitar''
atau ``kira-kira'' dalam pembahasan ini karena pernyataan tersebut
tidak sepenuhnya tepat, tetapi maksudnya sudah jelas. Terdapat galat
kecil ketika merepresentasikan $1/7$ sebagai bilangan riil titik-mengambang.
Hal yang sama berlaku bagi semua bilangan riil, kecuali suatu himpunan
berhingga.

Memang terdapat galat dalam representasi titik-mengambang bilangan
riil, tetapi galat itu selalu kecil dibandingkan dengan besar bilangan
riil yang direpresentasikan. Galat relatifnya sekitar $10^{-17}$,
sehingga pembahasan galat titik-mengambang mungkin tampak semata-mata
sebagai latihan akademis. Lagi pula, seberapa berartinya galat sebesar
$7.93(10)^{-18}$? Adakah yang akan kesal jika ia dijual sebuah cincin
selebar $.14285714285714284921$ inci padahal seharusnya lebarnya
$.14285714285714285714$ inci? Jelas tidak. Namun, yang perlu dikhawatirkan
bukan hanya galat dalam satu perhitungan (penjumlahan, pengurangan,
perkalian, atau pembagian). Metode numerik memerlukan puluhan, ribuan,
bahkan jutaan perhitungan. Galat-galat kecil dapat terakumulasi. Cobalah
percobaan berikut.

\subsubsection{Percobaan 1}

\noindent Gunakan kalkulator atau komputer Anda untuk menghitung bilangan-bilangan
$p_{0},p_{1},p_{2},\ldots,p_{7}$ sesuai aturan berikut.
\begin{itemize}
\item $p_{0}=\pi$
\item $p_{1}=10p_{0}-31$
\item $p_{2}=100p_{1}-41$
\item $p_{3}=100p_{2}-59$
\item $p_{4}=100p_{3}-26$
\item $p_{5}=100p_{4}-53$
\item $p_{6}=100p_{5}-58$
\item $p_{7}=100p_{6}-97$
\end{itemize}
Menurut kalkulator atau komputer Anda, $p_{7}$ kemungkinan menghasilkan
salah satu nilai berikut:

\begin{center}
\begin{tabular}{cc}
$0.93116$  & (Octave)\tabularnewline
$.9311599796346854$  & (Maxima)\tabularnewline
$1$ & (CASIO \textit{fx}-115ES)\tabularnewline
\end{tabular}
\par\end{center}

\noindent Namun, sedikit aljabar akan menunjukkan bahwa $p_{7}=10000000000000\pi-31415926535897$
secara eksak; hingga 6 tempat desimal, nilainya adalah $0.932384$.
Bandingkan nilai ini dengan bilangan yang Anda peroleh untuk $p_{7}$.
Meskipun $p_{0}$ merupakan hampiran yang sangat akurat terhadap $\pi$
(akurat hingga sekitar 16 tempat desimal), setelah hanya beberapa
perhitungan, galat pembulatan menyebabkan $p_{7}$ hanya memiliki
satu atau dua angka signifikan yang akurat!

Percobaan ini menyoroti penyebab terpenting galat titik-mengambang:
pengurangan dua bilangan yang hampir sama. Kita berulang kali mengurangkan
bilangan-bilangan yang angka puluhan dan satuannya sama. Dua angka
signifikan terdepan keduanya cocok. Sebagai contoh, $10\pi-31=31.415926\ldots-31$.
$10\pi$ dipertahankan akurat hingga sekitar 16 angka ($31.41592653589793$),
tetapi $10\pi-31$ dipertahankan akurat hanya hingga 14 angka signifikan
($0.41592653589793$). Setiap pengurangan berikutnya menurunkan akurasi
sebanyak dua angka signifikan lagi. Bahkan, $p_{7}$ direpresentasikan
hanya dengan 2 angka signifikan. Kita telah berulang kali mengurangkan
bilangan-bilangan yang hampir sama. Setiap kali melakukannya, sebagian
akurasi hilang. Galat titik-mengambang pun membesar.

Dalam perhitungan yang tidak melibatkan pengurangan dua besaran yang
hampir sama, galat algoritmik tetap perlu diperhatikan. Sebagai contoh,
misalkan $f(x)=\sin x$. Dari definisi turunan dapat dibuktikan bahwa
\[
f'(1)=\lim_{h\to0}\frac{\sin(1+h)-\sin(1-h)}{2h}.
\]
Karena itu, secara umum kita mengharapkan bahwa $\tilde{p}(h)=\frac{\sin(1+h)-\sin(1-h)}{2h}$
merupakan hampiran yang baik terhadap $f'(1)$ untuk nilai $h$ yang
kecil; dan semakin kecil $h$, semakin baik pula hampirannya. 

\subsubsection{Percobaan 2\label{subsec:derivativeExperiment}}

Dengan kalkulator atau komputer, hitung $\tilde{p}(h)$ untuk $h=10^{-2}$,
$h=10^{-3}$, dan seterusnya hingga $h=10^{-7}$. Hasil Anda semestinya
menyerupai tabel berikut:

\begin{center}
\begin{tabular}{c|c}
$h$ & $\tilde{p}(h)$\tabularnewline
\hline 
$10^{-2}$ & $0.5402933008747335$\tabularnewline
$10^{-3}$ & $0.5403022158176896$\tabularnewline
$10^{-4}$ & $0.5403023049677103$\tabularnewline
$10^{-5}$ & $0.5403023058569989$\tabularnewline
$10^{-6}$ & $0.5403023058958567$\tabularnewline
$10^{-7}$ & $0.5403023056738121$\tabularnewline
\end{tabular}
\par\end{center}

\noindent Kolom kedua dihitung menggunakan aritmetika hampiran (titik-mengambang),
sehingga secara teknis merupakan hampiran atas suatu hampiran. Karena
$f'(1)=\cos(1)\approx.5403023058681398$, setiap hampiran memang cukup
dekat dengan nilai eksak. Namun, jika kita mengamatinya lebih saksama,
masih ada hal lain yang perlu dikatakan. Pertama, galat algoritmik\index{galat!algoritmik}
dari $\tilde{p}(10^{-2})$ adalah 
\begin{eqnarray*}
|\tilde{p}^{*}(10^{-2})-f'(1)| & = & \left|50\left(\sin\left(\frac{101}{100}\right)-\sin\left(\frac{99}{100}\right)\right)-\cos(1)\right|\\
 & \approx & 9.00(10)^{-6}
\end{eqnarray*}
yang akurat hingga tiga angka signifikan. Artinya, jika kita menghitung
$\tilde{p}(10^{-2})$ dengan aritmetika eksak, yang di sini dinotasikan
dengan $\tilde{p}^{*}(10^{-2})$, nilainya tetap meleset dari $f'(1)$
sebesar kira-kira $9(10)^{-6}$. Galat titik-mengambang\index{galat!titik-mengambang}
hanya mengukur sejauh apa nilai terhitung dari $\tilde{p}(10^{-2})$
menyimpang dari nilai eksak $\tilde{p}^{*}(10^{-2})$ untuk $\tilde{p}(10^{-2})$.
Dengan kata lain, galat titik-mengambang diberikan oleh $\left|\tilde{p}^{*}-\tilde{p}\right|$:
\[
\left|50\left(\sin\left(\frac{101}{100}\right)-\sin\left(\frac{99}{100}\right)\right)-0.5402933008747335\right|\approx1.58(10)^{-17},
\]
yang sekecil mungkin seperti yang dapat diharapkan. Galat absolut
$|\tilde{p}(10^{-2})-f'(1)|=|0.5402933008747335-\cos(1)|$ pada dasarnya
sepenuhnya bersifat algoritmik. Galat pembulatan jauh lebih kecil
daripada galat algoritmik. Fakta bahwa kita menggunakan aritmetika
titik-mengambang dapat diabaikan. 

Sebaliknya, galat algoritmik dari $\tilde{p}(10^{-7})$ adalah 
\begin{eqnarray*}
|\tilde{p}^{*}(10^{-7})-f'(1)| & = & \left|5000000\left(\sin\left(\frac{10000001}{10000000}\right)-\sin\left(\frac{9999999}{10000000}\right)\right)-\cos(1)\right|\\
 & \approx & 9.00(10)^{-16}
\end{eqnarray*}
yang akurat hingga tiga angka signifikan. Namun, kita perlu sedikit
mengkhawatirkan galat titik-mengambang karena ${\displaystyle \sin\left(\frac{10000001}{10000000}\right)\approx0.8414710388}$
dan ${\displaystyle \sin\left(\frac{9999999}{10000000}\right)\approx.8414709307}$
hampir sama. Kita mengurangkan bilangan yang lima angka signifikan
terdepannya cocok! Galat titik-mengambangnya sekali lagi adalah $|\tilde{p}^{*}-\tilde{p}|$,
yaitu
\[
\left|5000000\left(\sin\left(\frac{10000001}{10000000}\right)-\sin\left(\frac{9999999}{10000000}\right)\right)-0.5403023056738121\right|\approx1.94(10)^{-10}.
\]
Galat ini mungkin tampak kecil, tetapi nilainya sangat besar dibandingkan
dengan galat algoritmik sekitar $9(10)^{-16}$. Jadi, dalam kasus
ini, pada dasarnya seluruh galat timbul karena kita menggunakan aritmetika
titik-mengambang! Kali ini, galat algoritmik jauh lebih kecil daripada
galat pembulatan. Untungnya, hal ini tidak sering terjadi sehingga
kita biasanya dapat memusatkan perhatian hanya pada galat algoritmik. 

\noindent\begin{digression}{Kekacauan} 

\noindent Edward Lorenz\index{Lorenz, Edward}, seorang meteorolog
di Massachusetts Institute of Technology, termasuk orang pertama yang
mengenali dan mempelajari fenomena matematis yang kini disebut kekacauan.
Pada awal tahun 1960-an, ia sedang berusaha memodelkan sistem cuaca
untuk memperbaiki peramalan cuaca. Menurut salah satu versi kisahnya,
ia ingin mengulangi perhitungan yang baru saja dilakukannya. Untuk
menghemat waktu, ia menggunakan kondisi awal yang sama seperti sebelumnya,
tetapi membulatkannya menjadi tiga angka signifikan alih-alih enam.
Dengan keyakinan penuh bahwa perhitungan baru akan menyerupai perhitungan
lama, ia pergi minum kopi lalu kembali untuk memeriksa hasilnya. Betapa
terkejutnya ia ketika melihat hasil yang sama sekali berbeda! Ia mengulangi
prosedur itu beberapa kali dan setiap kali mendapati bahwa variasi
awal yang kecil menghasilkan variasi jangka panjang yang besar. Apakah
ini sekadar kasus galat titik-mengambang? Bukan. Berikut versi yang
cukup disederhanakan dari apa yang terjadi. Misalkan $f(x)=4x(1-x)$
dan tetapkan $p_{0}=1/7$. Sekarang hitung $p_{1}=f(p_{0})$, $p_{2}=f(p_{1})$,
$p_{3}=f(p_{2})$, dan seterusnya hingga diperoleh $p_{40}=f(p_{39})$.
Anda seharusnya memperoleh $p_{40}\approx0.080685$. Sekarang tetapkan
$p_{0}=1/7+10^{-12}$ (agar kita dapat menjalankan perhitungan yang
sama, tetapi dengan nilai awal yang berbeda dari nilai semula hanya
sebesar $10^{-12}$). Hitung seperti sebelumnya: $p_{1}=f(p_{0})$,
$p_{2}=f(p_{1})$, $p_{3}=f(p_{2})$, dan seterusnya hingga diperoleh
$p_{40}=f(p_{39})$. Kali ini Anda seharusnya memperoleh $p_{40}\approx0.91909$---hasil
yang sama sekali berbeda! Jika Anda mengulangi kedua perhitungan dengan
aritmetika 100 angka signifikan, Anda akan mendapati bahwa memulai
dengan $p_{0}=1/7$ menghasilkan $p_{40}\approx.080736$, sedangkan
memulai dengan $p_{0}=1/7+10^{-12}$ menghasilkan $p_{40}\approx0.91912$.
Dengan kata lain, bukan penggunaan hampiran titik-mengambang yang
membuat kedua perhitungan tersebut memberikan hasil yang sangat berbeda.
Penggunaan aritmetika 1000 angka signifikan tidak akan mengubah kesimpulan,
demikian pula perhitungan yang lebih presisi apa pun. Ini merupakan
demonstrasi dari apa yang disebut kepekaan terhadap kondisi awal,
suatu ciri semua sistem kacau, termasuk cuaca. Variasi yang sangat
kecil pada suatu saat menghasilkan variasi yang sangat besar kemudian.
Dan ``galat'' tersebut bersifat algoritmik. Inilah prinsip dasar
yang membuat peramalan cuaca jangka panjang mustahil. Dalam kata-kata
Edward Lorenz, ``Mengingat ketidakakuratan dan ketidaklengkapan pengamatan
cuaca yang tidak dapat dihindari, peramalan jangka sangat panjang
yang presisi tampaknya tidak mungkin ada.''
\begin{description}
\item [{Referensi}] \cite{lorenz,gulick,mirror}
\end{description}
\noindent\end{digression} 

\subsubsection{Percobaan 3}

Misalkan $a=77617$ dan $b=33096$, lalu hitung
\[
333.75b^{6}+a^{2}(11a^{2}b^{2}-b^{6}-121b^{4}-2)+5.5b^{8}+\frac{a}{2b}.
\]
Anda mungkin memperoleh bilangan seperti $-1.180591620717411(10)^{21}$
meskipun nilai eksaknya adalah
\[
-\frac{54767}{66192}\approx-.8273960599468214.
\]
Galatnya luar biasa besar! Namun, hal ini bukan karena kalkulator
atau komputer Anda bermasalah ketika menghitung setiap suku dengan
tingkat akurasi yang wajar. Cobalah sendiri.
\begin{eqnarray*}
333.75b^{6} & = & 438605750846393161930703831040\\
a^{2}(11a^{2}b^{2}-b^{6}-121b^{4}-2) & = & -7917111779274712207494296632228773890\\
5.5b^{8} & = & 7917111340668961361101134701524942848\\
\frac{a}{2b} & = & \frac{77617}{66192}\approx1.172603940053179
\end{eqnarray*}
Penyebab buruknya hasil perhitungan adalah pengurangan nilai-nilai
yang hampir sama setelah setiap suku dihitung. $a^{2}(11a^{2}b^{2}-b^{6}-121b^{4}-2)$
dan $5.5b^{8}$ memiliki tanda berlawanan dan cocok pada 7 angka signifikan
terbesar, sehingga menghitung jumlahnya menurunkan akurasi sekitar
7 angka signifikan. Lebih buruk lagi, $a^{2}(11a^{2}b^{2}-b^{6}-121b^{4}-2)+5.5b^{8}=-438605750846393161930703831042$,
yang bertanda berlawanan dengan $333.75b^{6}$ dan cocok dengannya
pada setiap nilai tempat kecuali satuan. Itu berarti 29 angka! Jadi,
kita kehilangan lagi 29 angka signifikan akurasi ketika menambahkan
jumlah ini pada $333.75b^{6}$. Jika dihitung secara eksak, jumlah
$333.75b^{6}+a^{2}(11a^{2}b^{2}-b^{6}-121b^{4}-2)+5.5b^{8}$ adalah
$-2$. Namun, perhitungan harus dilakukan hingga 37 angka signifikan
agar hasil ini terlihat. Perhitungan yang hanya menggunakan sekitar
16 angka signifikan, seperti yang dilakukan kebanyakan kalkulator
dan komputer, menghasilkan 0 angka signifikan yang akurat karena 36
angka akurasi hilang selama perhitungan. Itulah sebabnya Anda dapat
memperoleh bilangan seperti $-1.180591620717411(10)^{21}$ sebagai
jawaban akhir alih-alih jawaban eksak $\frac{a}{2b}-2\approx-.8273960599468214$.

Yang mungkin lebih mengejutkan adalah bahwa penyusunan ulang sederhana
atas ekspresi tersebut menghasilkan nilai yang sama sekali berbeda.
Cobalah menghitung 
\[
(333.75-a^{2})b^{6}+a^{2}(11a^{2}b^{2}-121b^{4}-2)+5.5b^{8}+\frac{a}{2b}
\]
sebagai gantinya. Kali ini Anda mungkin memperoleh bilangan seperti
$1.172603940053179$. Sekali lagi hasilnya sama sekali tidak akurat,
dan penyebabnya sama. Kali ini setiap sukunya adalah
\begin{eqnarray*}
(333.75-a^{2})b^{6} & = & -7917110903377385049079188237280149504\\
a^{2}(11a^{2}b^{2}-121b^{4}-2) & = & -437291576312021946464244793346\\
5.5b^{8} & = & 7917111340668961361101134701524942848\\
\frac{a}{2b} & = & \frac{77617}{66192}\approx1.172603940053179
\end{eqnarray*}
sehingga masalahnya tetap ada. Kita masih harus mengurangkan bilangan-bilangan
yang nilainya hampir sama. Perbedaan antara perhitungan ini dan perhitungan
sebelumnya terletak pada pembulatan. Dalam kasus pertama, pembulatan
membuat dua bilangan besar berbeda pada angka signifikan terakhirnya,
sehingga jumlahnya menjadi sangat besar. Dalam kasus kedua, jumlah
tiga suku pertama ternyata $0$ karena bilangan-bilangan besar tersebut
cocok pada semua angka signifikan. Perhatikan bahwa dalam kasus kedua,
hasil akhirnya hanyalah nilai $\frac{a}{2b}$.

Seperti diperlihatkan oleh contoh-contoh ini, kadang galat titik-mengambang
dan kadang galat algoritmik dapat merusak suatu perhitungan. Namun,
secara umum galat titik-mengambang sangat sulit dideteksi. Galat algoritmik
jauh lebih mudah dianalisis. Selain itu, sebagian besar algoritma
yang akan kita pelajari tidak rentan terhadap galat titik-mengambang.
Dalam hampir semua kasus, bagian terbesar galat akan bersifat algoritmik.
\begin{description}
\item [{Referensi}] \cite{rump,loh}
\end{description}

\subsection{Konsep Utama}
\begin{description}
\item [{$p$}] Nilai eksak yang sedang dihampiri.
\item [{$\tilde{p}$}] Suatu hampiran terhadap nilai $p$.
\item [{$\tilde{p}^{*}$}] Suatu hampiran terhadap nilai $p$ yang dihitung
menggunakan aritmetika eksak.
\item [{Galat~absolut\index{galat!absolut}:}] $|\tilde{p}-p|$ disebut
galat absolut ketika menggunakan $\tilde{p}$ untuk menghampiri nilai
$p$.
\item [{Galat~relatif\index{galat!relatif}:}] ${\displaystyle \frac{|\tilde{p}-p|}{|p|}}$
disebut galat relatif ketika menggunakan $\tilde{p}$ untuk menghampiri
nilai $p$.
\item [{Akurasi\index{akurasi}:}] Kita mengatakan bahwa $\tilde{p}$ akurat
hingga $n$ angka signifikan jika $n$ angka signifikan terdepan dari
$\tilde{p}$ cocok dengan angka-angka dari $p$. Secara lebih tepat,
kita mengatakan bahwa $\tilde{p}$ akurat hingga $d(\tilde{p})=\log\left|\frac{p}{\tilde{p}-p}\right|$
angka signifikan.
\item [{Aritmetika~titik-mengambang\index{aritmetika titik-mengambang}:}] Aritmetika
yang menggunakan bilangan-bilangan yang direpresentasikan dengan banyak
angka signifikan yang tetap.
\item [{Galat~algoritmik\index{galat!algoritmik}:}] Galat yang semata-mata
disebabkan oleh algoritma atau persamaan yang digunakan dalam hampiran,
$|\tilde{p}^{*}-p|$, dengan $\tilde{p}^{*}$ merupakan hampiran terhadap
$p$ yang dihitung menggunakan aritmetika eksak.
\item [{Galat~pemenggalan\index{galat!pemenggalan}:}] Galat algoritmik
akibat penggunaan jumlah parsial sebagai pengganti suatu deret. Pada
jenis galat ini, ekor deret dipenggal---dari sinilah namanya berasal.
\item [{Galat~titik-mengambang\index{galat!titik-mengambang}:}] Galat
yang semata-mata disebabkan oleh perhitungan menggunakan aritmetika
titik-mengambang, $|\tilde{p}^{*}-\tilde{p}|$ dengan $\tilde{p}$
dihitung menggunakan aritmetika titik-mengambang, $\tilde{p}^{*}$
dihitung menggunakan aritmetika eksak, dan keduanya dihitung menurut
rumus atau algoritma yang sama.
\item [{Galat~pembulatan\index{galat!pembulatan}:}] Nama lain untuk galat
titik-mengambang.
\end{description}

\subsection{Octave}

Perhitungan dalam bagian ini dapat dilakukan dengan mudah menggunakan
Octave. Yang Anda perlukan hanyalah operasi aritmetika\index{Octave!operasi aritmetika}
dan beberapa fungsi standar\index{Octave!fungsi standar} seperti
nilai absolut, sinus, dan kosinus. Untungnya, semua ini cukup mudah
dilakukan dengan Octave. Operasi aritmetika dilakukan hampir sama
seperti pada kalkulator. Hanya ada satu perbedaan penting. Kebanyakan
kalkulator akan menerima ekspresi seperti \texttt{3x} dan memahami
bahwa yang Anda maksud adalah $3\times x$, tetapi Octave tidak demikian.
Ekspresi \texttt{3x} menyebabkan galat sintaks di Octave. Octave mengharuskan
Anda menyatakan operasinya secara eksplisit, seperti dalam \texttt{3{*}x}. 

Fungsi standar seperti nilai absolut, sinus, dan kosinus (serta banyak
fungsi lain) memiliki singkatan sederhana di Octave. Semua fungsi
tersebut menerima satu argumen, atau masukan. Ingatlah notasi fungsi;
dengan demikian akan jelas cara mencari sinus atau nilai absolut suatu
bilangan. Anda perlu mengetikkan nama fungsi, tanda kurung buka, argumen,
dan tanda kurung tutup, seperti dalam \texttt{sin(7.2)}. Beberapa
fungsi umum beserta singkatannya tercantum dalam Tabel \ref{tab:commonOctave}.
\begin{table}[h]
\hrule \caption{\label{tab:commonOctave}Beberapa fungsi umum dan singkatannya dalam
Octave.}

\medskip{}

\centering{}%
\begin{tabular}{ll||ll||ll}
Fungsi & Octave & Fungsi & Octave & Fungsi & Octave\tabularnewline
\hline 
$n!$ & \texttt{factorial(n)} & $\sin(x)$ & \texttt{sin(x)} & $\cos(x)$ & \texttt{cos(x)}\tabularnewline
$|x|$ & \texttt{abs(x)} & $\tan(x)$ & \texttt{tan(x)} & $\cot(x)$ & \texttt{cot(x)}\tabularnewline
$e^{x}$ & \texttt{exp(x)} atau \texttt{e\textasciicircum x} & $\sec(x)$ & \texttt{sec(x)} & $\csc(x)$ & \texttt{csc(x)}\tabularnewline
$\ln(x)$ & \texttt{log(x)} & $\sin^{-1}(x)$ & \texttt{asin(x)} & $\cos^{-1}(x)$ & \texttt{asin(x)}\tabularnewline
$\log(x)$ & \texttt{log10(x)} & $\tan^{-1}(x)$ & \texttt{atan(x)} & $\cot^{-1}(x)$ & \texttt{acot(x)}\tabularnewline
$\sqrt{x}$ & \texttt{sqrt(x)} & $\sqrt[3]{x}$ & \texttt{cbrt(x)} & $\sqrt[n]{x}$ & \texttt{nthroot(x)}\tabularnewline
$b^{x}$ & \texttt{b\textasciicircum x} & $\sinh(x)$ & \texttt{sinh(x)} & $\cosh(x)$ & \texttt{cosh(x)}\tabularnewline
$\lfloor x\rfloor$ & \texttt{floor(x)} & $\lceil x\rceil$ & \texttt{ceil(x)} &  & \tabularnewline
\end{tabular}\medskip{}
\hrule 
\end{table}
 Fungsi dan operasi aritmetika dapat digabungkan dengan cara yang
wajar. Beberapa contoh dari bagian ini ditampilkan dalam Tabel \ref{tab:errorExamples}.
\begin{table}[h]
\hrule \caption{\label{tab:errorExamples}Perhitungan beberapa ekspresi dengan Octave.}

\medskip{}

\centering{}%
\begin{tabular}{lll}
Ekspresi & Octave & Hasil\tabularnewline
\hline 
$|\frac{22}{7}-\pi|$ & \texttt{abs(22/7-pi)} & \texttt{0.0012645}\tabularnewline
$\frac{|\frac{16525}{54}-\pi^{5}|}{|\pi^{5}|}$ & \texttt{abs(16525/54-pi\textasciicircum 5)/abs(pi\textasciicircum 5)} & \texttt{3.8111e-06}\tabularnewline
$\frac{\sin(1.01)-\sin(0.99)}{0.02}$ & \texttt{(sin(1.01)-sin(0.99))/0.02} & \texttt{0.54029}\tabularnewline
\end{tabular}\medskip{}
\hrule 
\end{table}
Ada dua hal yang perlu diperhatikan. Pertama, notasi Octave sangat
mirip dengan notasi kalkulator. Kedua, secara bawaan Octave menampilkan
hasil menggunakan 5 angka signifikan. Namun, jangan mengira Octave
hanya menghitung lima angka hasil tersebut. Sebenarnya, Octave telah
menghitung sedikitnya 15 angka dengan benar. Jika Anda ingin melihat
angka-angka itu, gunakan perintah \texttt{format('long')} \index{Octave!format@\texttt{format}}.
Perintah ini hanya perlu digunakan sekali dalam setiap sesi. Semua
bilangan yang dicetak setelah perintah ini dijalankan akan ditampilkan
dengan 15 angka signifikan. Sebagai contoh, \texttt{1/7} akan menghasilkan
\texttt{0.142857142857143} alih-alih hanya \texttt{0.14286}. Jika
ingin kembali ke format bawaan, gunakan perintah \texttt{format()}
tanpa argumen. Kita akan membahas pengendalian keluaran yang lebih
terperinci nanti. Untuk sekarang, berikut beberapa cara melakukan
Percobaan 1 menggunakan Octave. Perbedaannya hanya pada banyaknya
keluaran dan format keluarannya. Bilangan-bilangannya dihitung dengan
cara dan presisi yang tepat sama.

\subsubsection{Percobaan 1 di Octave, contoh 1}

\begin{verbatim}
octave:1> p0=pi;
octave:2> p1=10*p0-31; p2=100*p1-41; p3=100*p2-59;
octave:3> p4=100*p3-26; p5=100*p4-53; p6=100*p5-58;
octave:4> p7=100*p6-97
p7 =  0.93116 
\end{verbatim}

\subsubsection{Percobaan 1 di Octave, contoh 2}

\begin{verbatim}
octave:1> format('long')
octave:2> p0=pi
p0 =  3.14159265358979
octave:3> p1=10*p0-31
p1 =  0.415926535897931
octave:4> p2=100*p1-41
p2 =  0.592653589793116
octave:5> p3=100*p2-59
p3 =  0.265358979311600
octave:6> p4=100*p3-26
p4 =  0.535897931159980
octave:7> p5=100*p4-53
p5 =  0.589793115997963
octave:8> p6=100*p5-58
p6 =  0.979311599796347
octave:9> p7=100*p6-97
p7 =  0.931159979634685 
\end{verbatim}

\subsubsection{Percobaan 1 di Octave, contoh 3}

\begin{verbatim}
octave:1> 10*pi-31
ans =  0.41593
octave:2> 100*ans-41
ans =  0.59265
octave:3> 100*ans-59
ans =  0.26536
octave:4> 100*ans-26
ans =  0.53590
octave:5> 100*ans-53
ans =  0.58979
octave:6> 100*ans-58
ans =  0.97931
octave:7> 100*ans-97
ans =  0.93116 
\end{verbatim}

\subsubsection{Percobaan 3 di Octave}

\begin{verbatim}
octave:1> a=77617;
octave:2> b=33096;
octave:3> t1=333.75*b^6;
octave:4> t2=a^2*(11*a^2*b^2-b^6-121*b^4-2);
octave:5> t3=5.5*b^8;
octave:6> t4=a/(2*b);
octave:7> t1+t2+t3+t4
ans = -1.18059162071741e+21
octave:8> t1=(333.75-a^2)*b^6;
octave:9> t2=a^2*(11*a^2*b^2-121*b^4-2);
octave:10> t1+t2+t3+t4
ans =  1.17260394005318 
\end{verbatim}

\noindent Pada akhirnya, cara yang Anda pilih untuk mengerjakan suatu
latihan di Octave bergantung pada preferensi dan tujuan Anda. Ajukanlah
pertanyaan seperti berikut kepada diri sendiri. Berapa angka signifikan
yang saya perlukan? Berapa banyak hasil antara yang perlu saya lihat?
Hasil yang mana? Jawaban atas pertanyaan semacam itu seharusnya memandu
penyelesaian Anda.

Jika diperlukan, Octave memiliki singkatan untuk sebagian besar konstanta
umum. Tabel \ref{tab:constants} menampilkan tiga konstanta yang paling
umum.
\begin{table}[h]
\hrule \caption{\label{tab:constants}Beberapa konstanta \index{Octave!konstanta}
dalam Octave.}

\medskip{}

\centering{}%
\begin{tabular}{lll}
Konstanta & Octave & Hasil\tabularnewline
\hline 
$e$ & \texttt{e} & \texttt{2.7183}\tabularnewline
$\pi$ & \texttt{pi} & \texttt{3.1416}\tabularnewline
$i$ & \texttt{i atau j} & \texttt{0 + 1i}\tabularnewline
\end{tabular}\medskip{}
\hrule 
\end{table}

\noindent\startexercises

\subsection{Latihan}
\begin{enumerate}
\item Selain galat pembulatan, apa yang dapat berdampak buruk terhadap akurasi
suatu perhitungan numerik?
\item Hitung galat absolut dan galat relatif ketika menghampiri $\pi$ dengan
3.
\item \label{ex:absoluteError}Hitung galat absolut ketika menghampiri $p$
dengan $\tilde{p}$.

\begin{enumerate}
\item \label{exc:abserr1}$p=123$;~~ $\tilde{p}=\frac{1106}{9}$ \hasasolution 
\item $p=\frac{1}{e}$;~~ $\tilde{p}=.3666$
\item \label{exc:abserr2}$p=2^{10}$;~~ $\tilde{p}=1000$ \hasasolution 
\item $p=24$;~~ $\tilde{p}=48$
\item \octave\ \label{exc:abserr3}$p=\pi^{-7}$;~~ $\tilde{p}=10^{-4}$ \hasasolution 
\item \octave\ $p=(0.062847)(0.069234)$;~~ $\tilde{p}=0.0042$
\end{enumerate}
\item \label{exc:abserr15}Hitung galat relatif untuk hampiran-hampiran
pada soal \ref{ex:absoluteError}. \haspartwithsolution 
\item \label{exc:abserr16}Berapa banyak angka signifikan yang akurat pada
hampiran-hampiran dalam soal \ref{ex:absoluteError} tersebut? \haspartwithsolution 
\item Hitung galat absolut dan galat relatif ketika menghampiri $p$ dengan
$\tilde{p}$.

\begin{enumerate}
\item $p=\sqrt{2}$, $\tilde{p}=1.414$
\item $p=10^{\pi}$, $\tilde{p}=1400$
\item $p=9!$, $\tilde{p}=\sqrt{18\pi}(9/e)^{9}$
\end{enumerate}
\item \octave\ \label{enu:one-over-pi}Hitung ${\displaystyle \frac{1103\sqrt{8}}{9801}}$
menggunakan Octave.
\item \octave\ Bilangan pada soal \ref{enu:one-over-pi} merupakan hampiran
terhadap $1/\pi$. Dengan Octave, tentukan galat absolut dan galat
relatif hampiran tersebut.
\item \octave\ Dengan Octave, hitung

\begin{enumerate}
\item $\lfloor\ln(234567)\rfloor$
\item $e^{\lceil\ln(234567)\rceil}$
\item $\sqrt[3]{\lfloor\sin(e^{5.2})\rfloor}$
\item $-e^{i\pi}$
\item $4\tan^{-1}(1)$
\item ${\displaystyle \frac{\lfloor\cos(3)-\sqrt[5]{\ln(3)}\rfloor}{\lceil\arctan(3)-e^{3}\rceil}}$
\end{enumerate}
\item \octave\ Tentukan $f(2)$ menggunakan Octave.

\begin{enumerate}
\item \label{exc:abserr12}$f(x)=e^{\sin(x)}$ \hasasolution 
\item $f(x)=\sin\left(e^{x}\right)$
\item \label{exc:abserr13}${\displaystyle f(x)=\tan^{-1}(x-0.429)}$ \hasasolution 
\item $f(x)=x-\tan^{-1}(0.429)$
\item \label{exc:abserr14}$f(x)=10^{x}/5!$ \hasananswer 
\item $f(x)=5!/x^{10}$
\end{enumerate}
\item \octave\ Semua persamaan berikut benar secara matematis. Meskipun
demikian, galat titik-mengambang membuat beberapa di antaranya dianggap
salah oleh Octave. Yang mana saja? PETUNJUK: Gunakan operator Boolean
\texttt{==} untuk memeriksanya. Sebagai contoh, untuk memeriksa apakah
$\sin(0)=0$, ketik \texttt{sin(0)==0} di Octave. \texttt{ans=1} berarti
benar (kedua ruas sama menurut Octave---tidak ada galat pembulatan),
sedangkan \texttt{ans=0} berarti salah (kedua ruas tidak sama menurut
Octave---terdapat galat pembulatan).

\begin{enumerate}
\item $(2)(12)=9^{2}-4(9)-21$
\item $e^{3\ln(2)}=8$
\item $\ln(10)=\ln(5)+\ln(2)$
\item $g(\frac{1+\sqrt{5}}{2})=\frac{1+\sqrt{5}}{2}$ dengan $g(x)=\sqrt[3]{x^{2}+x}$
\item $\lfloor153465/3\rfloor=153465/3$
\item $3\pi^{3}+7\pi^{2}-2\pi+8=((3\pi+7)\pi-2)\pi+8$
\end{enumerate}
\item \label{ex:absoluteError2}Tentukan suatu hampiran $\tilde{p}$ terhadap
$p$ dengan galat absolut $.001$.

\begin{enumerate}
\item \label{exc:abserr4}$p=\pi$ \hasasolution 
\item $p=\sqrt{5}$
\item \label{exc:abserr5}$p=\ln(3)$ \hasasolution 
\item $p=\sqrt{23}^{\sqrt{23}}$
\item \label{exc:abserr6}$p=\frac{10}{\ln(1.1)}$ \hasasolution 
\item $p=\tan(1.57079)$
\end{enumerate}
\item \label{exc:abserr17}Tentukan suatu hampiran $\tilde{p}$ terhadap
$p$ dengan galat relatif $.001$ untuk setiap nilai $p$ pada soal
\ref{ex:absoluteError2}. \haspartwithsolution 
\item \label{ex:absoluteError3}$\tilde{p}$ menghampiri nilai apa dengan
galat absolut $.0005$?

\begin{enumerate}
\item \label{exc:abserr7}$\tilde{p}=.2348263818643$ \hasananswer 
\item $\tilde{p}=23.89627345677$
\item $\tilde{p}=-8.76257664363$
\end{enumerate}
\item \label{exc:abserr18}Ulangi soal \ref{ex:absoluteError3}, tetapi
gunakan galat relatif $.0005$. \haspartwithanswer 
\item \label{exc:abserr8}$\tilde{p}$ menghampiri $p$ dengan galat absolut
$\frac{1}{100}$ dan galat relatif $\frac{3}{100}$. Tentukan $p$
dan $\tilde{p}$. \hasananswer 
\item $\tilde{p}$ menghampiri $p$ dengan galat absolut $\frac{3}{100}$
dan galat relatif $\frac{1}{100}$. Tentukan $p$ dan $\tilde{p}$.
\item Misalkan $\tilde{p}$ harus menghampiri $p$ dengan galat relatif
paling besar $10^{-3}$. Tentukan interval terbesar tempat $\tilde{p}$
harus berada jika $p=900$.
\item Bilangan $e$ dapat didefinisikan oleh $e=\sum_{n=0}^{\infty}(1/n!)$.
Hitung galat absolut dan galat relatif pada hampiran-hampiran berikut
terhadap $e$:

\begin{enumerate}
\item ${\displaystyle \sum_{n=0}^{5}\frac{1}{n!}}$
\item ${\displaystyle \sum_{n=0}^{10}\frac{1}{n!}}$
\end{enumerate}
\item Rasio emas, ${\displaystyle \frac{1+\sqrt{5}}{2}}$, ditemukan di
berbagai tempat dalam alam maupun matematika. Sebagai contoh, jika
$F_{n}$ adalah bilangan Fibonacci ke-$n$, maka
\[
\lim_{n\to\infty}\frac{F_{n+1}}{F_{n}}=\frac{1+\sqrt{5}}{2}
\]
Karena itu, $F_{11}/F_{10}$ dapat digunakan sebagai hampiran terhadap
rasio emas. Tentukan galat relatif hampiran ini. PETUNJUK: Barisan
Fibonacci didefinisikan oleh $F_{0}=1$, $F_{1}=1$, $F_{n}=F_{n-1}+F_{n-2}$
untuk $n\geq2$.
\item \label{exc:abserr9}Tentukan nilai $p$ dan $\tilde{p}$ sedemikian
sehingga galat relatif sama dengan galat absolut. Nyatakan secara
umum kondisi yang menyebabkan hal ini terjadi. \hasananswer 
\item Tentukan nilai $p$ dan $\tilde{p}$ sedemikian sehingga galat relatif
lebih besar daripada galat absolut. Nyatakan secara umum kondisi yang
menyebabkan hal ini terjadi.
\item Tentukan nilai $p$ dan $\tilde{p}$ sedemikian sehingga galat relatif
lebih kecil daripada galat absolut. Nyatakan secara umum kondisi yang
menyebabkan hal ini terjadi.
\item Hitung (i) $\tilde{p}$ dengan kalkulator atau komputer (perhitungan
titik-mengambang), (ii) galat absolut, $|\tilde{p}-p|$, dan (iii)
galat relatif, $\frac{|\tilde{p}-p|}{|p|}$. Kemudian gunakan nilai
$\tilde{p}^{*}$ yang diberikan untuk menghitung (iv) galat algoritmik,
$|\tilde{p}^{*}-p|$ dan (v) galat pembulatan, $|\tilde{p}^{*}-\tilde{p}|$.

\begin{enumerate}
\item \label{exc:abserr10}Misalkan $f(x)=x^{4}+7x^{3}-63x^{2}-295x+350$
dan misalkan $p=f'(-2)$. Nilai $\tilde{p}=\frac{f(-2+10^{-7})-f(-2-10^{-7})}{2(10)^{-7}}$
merupakan hampiran yang baik terhadap $p$. $\tilde{p}$ tepat sama
dengan $8.99999999999999$. \hasananswer 
\item Misalkan $f'(x)=e^{x}\sin(10x)$ dan $f(0)=0$, serta misalkan $p=f(1)$.
Dapat ditunjukkan bahwa $p=\frac{1}{101}e(\sin10-10\cos10)+\frac{10}{101}$.
Metode Euler menghasilkan hampiran $\tilde{p}=\frac{1}{10}\sum_{i=1}^{9}e^{i/10}\sin i$.
Hingga 28 angka signifikan yang akurat, $\tilde{p}$ adalah $0.3550449718187918680449850763$.
Artinya, gunakan $\tilde{p}^{*}=0.3550449718187918680449850763$.
\item Misalkan $a_{0}=\frac{5+\sqrt{5}}{8}$ dan $a_{n+1}=4a_{n}(1-a_{n})$,
lalu perhatikan $p=a_{51}$. Dapat ditunjukkan bahwa $p=a_{51}=\frac{5-\sqrt{5}}{8}$.
Algoritma paling langsung untuk menghitung $a_{51}$ adalah menghitung
$a_{1},a_{2},a_{3},\ldots a_{51}$ secara berurutan menurut relasi
rekurensi yang diberikan. Gunakan algoritma ini untuk menghitung $\tilde{p}^{*}$
dan $\tilde{p}$.
\end{enumerate}
\end{enumerate}
\finishexercises 
